\section{Approach}
\label{sec:approach}

\subsection{Overview}

MPTA-Repair is a counterexample-guided framework designed for the
multi-property clock-bound repair of timed automata models. The
framework takes as input a faulty model $M$, a property set
$\Phi=\{\varphi_1,\ldots,\varphi_m\}$, and configuration parameters such
as a global timeout. Instead of requiring users to manually specify
repair variables, MPTA-Repair extracts repairable bounds from numerical
clock constraints in transition guards and location invariants. The
output is either a repaired model $M'$ along with a set of clock-bound
edits, or a failure report indicating that no admissible full-property
repair was identified within the given time budget.

MPTA-Repair is based on two observations. First, while candidate repair
generation can be driven by local diagnostic traces, candidate
validation must remain global. A candidate that repairs one violated
property might induce failures in others, because shared clocks,
locations, transitions, or synchronizations often participate in
multiple timing requirements. Consequently, MPTA-Repair iteratively
collects counterexamples from subsequent violations to progressively
constrain the space of admissible repair candidates. This refinement
loop continues until a candidate satisfies the entire property set and
passes admissibility verification, or the given time budget is
exhausted.

Second, different properties and diagnostic traces are often related
through shared repairable bounds. A single erroneous timing bound can
contribute to multiple violations, and distinct traces may correspond to
different executions of the same underlying fault. To leverage this
interdependence, MPTA-Repair maintains a joint pool of property-trace
obligations and searches for a clock-bound assignment that
simultaneously satisfies all encoded constraints while minimizing the
total modification cost. Relational analysis across properties, traces,
and repairable bounds serves to minimize redundant computations and
guide the search space exploration, though the generated candidates must
ultimately pass verification and admissibility checks.

In operation, MPTA-Repair initially verifies the model against the
property set and collects timed diagnostic traces (TDTs) for any
violations. It then extracts repairable clock bounds and constructs a
joint MaxSMT repair problem from the active property-trace pool. The
computed assignment is instantiated into a candidate model, which
undergoes validation against the complete property set and a
functional-equivalence admissibility criterion. If validation fails,
newly identified diagnostic traces are integrated into the pool, and the
refinement loop repeats.

Figure~\ref{fig:mpta-repair-workflow} illustrates the overall workflow.
The initial verification stage populates the property-trace pool.
Subsequently, the relation-guided optimization phase structures the
dependencies among properties, counterexamples, and repairable bounds.
Finally, the MaxSMT solving stage generates candidate edits. The
validation component either accepts an admissible repair or feeds new
diagnostic evidence back into the subsequent refinement round.

\begin{figure*}[t]
  \centering
  \includegraphics[width=0.98\textwidth]{mpta_repair_framework.pdf}
  \caption{Overview of the MPTA-Repair workflow.}
  \label{fig:mpta-repair-workflow}
\end{figure*}

\subsection{Joint Trace Encoding and Optimization Objective}

MPTA-Repair adopts the bound-variation paradigm for clock-bound repair
\cite{koelbl2019clock}. Each modifiable numerical bound in a transition
guard or location invariant is associated with a variation variable,
ensuring that the repaired constraint retains the original clock
reference and comparison operator while incorporating a shifted
numerical value. Consequently, the repair space is restricted to
clock-bound constants, whereas transition structures, location
topologies, synchronization labels, reset sets, and discrete updates
remain unaltered.

For each diagnostic trace, the framework constructs a TARTAR-style trace
constraint system \cite{koelbl2019clock} encompassing standard
conditions for clock initialization, time elapse, clock resets, guard
satisfaction, and invariant satisfaction. By replacing original clock
bounds with their varied counterparts, this formulation preserves the
discrete trace skeleton within the solver while allowing its timing
region to change. For a violated property and diagnostic trace, the
resulting obligation serves two purposes. It ensures post-repair
executability of the discrete trace skeleton to prevent trivial
elimination of functional paths, and it requires that no feasible timing
realization of this skeleton violates the target property. Local
obligations therefore isolate violating timed behaviors without blocking
the underlying functional execution.

To guarantee physical validity during candidate generation, static
consistency constraints are embedded directly into the formulation
\cite{demoura2008z3}. Repaired bounds must lie within the supported
numerical domain. Furthermore, when a lower bound and an upper bound
jointly constrain the same clock within a local region, their repaired
values must maintain a consistent timing interval.

From the active pool of property-trace obligations, MPTA-Repair
assembles a joint MaxSMT problem. The hard constraints include all
encoded trace obligations, static consistency requirements,
well-formedness conditions for repaired bounds, and blocking constraints
derived from previously rejected candidates. The optimization objective
is governed by a hierarchical set of simple soft constraints
\cite{sebastiani2015optimization}. The primary objective minimizes the
number of altered clock bounds by maximizing soft constraints of the
form $v=0$. Among candidate solutions with the same number of
modifications, a secondary objective minimizes the total magnitude of
the variations. This optimization strategy aligns with the engineering
goal of generating minimal, human-inspectable repairs and deliberately
avoids speculative metrics such as semantic risk or causal relevance.
Relation analysis is instead used externally to select or rank repair
variables and obligations before and around the solving phase.

\subsection{Relation-Guided Search Optimization}

During successive refinement rounds, the accumulation of violated
properties, diagnostic traces, and editable clock bounds can
significantly expand the search space \cite{clarke2000cegar}.
Unstructured encoding of all property-trace obligations and repairable
bounds increases the complexity of the MaxSMT problem and can induce
redundant solving cycles over overlapping evidence. To mitigate this
overhead, MPTA-Repair introduces relation-guided optimization to prune,
group, and rank variables and constraints. These relation-based
heuristics serve strictly as search optimizations; validity and
admissibility of generated candidates are still guaranteed through
full-property regression and formal verification.

At the property layer, MPTA-Repair operates on safety properties
normalized into an $\mathtt{A}[]$ fragment over location predicates and
clock or integer constraints \cite{vogel2022patterns}. Relation analysis
uses conservative syntactic and structural rules to evaluate equivalence
and dominance. Two properties are equivalent only when their normalized
representations are syntactically identical. Dominance is established
only when clear logical subsumption holds, for example when two
upper-bound properties share identical structures but differ in their
numerical thresholds and one threshold strictly implies the other.
Shared clocks, target locations, and template structures are tracked as
dependency metadata, but they do not independently justify removing
obligations.

At the counterexample layer, the framework maintains a managed pool of
TDTs indexed by their corresponding violated properties. Redundant
traces are eliminated through exact-duplicate removal. The remaining
TDTs are annotated with structural features, including traversed
transitions, visited locations, active clocks, guard and invariant
bounds, and originating properties \cite{koelbl2019clock}. Structural
similarity, common prefixes, and shared repair variables are used to
cluster related traces and orchestrate joint solving. A trace obligation
is deactivated only when an exact duplicate is detected or when a sound
subsumption check proves that its local constraint profile is covered by
an existing encoded obligation.

At the repair-bound layer, MPTA-Repair establishes a mapping between
each active property-trace obligation and the clock bounds capable of
influencing its feasibility. This mapping restricts the active decision
variables in the MaxSMT problem to those directly relevant to the
current counterexample pool, while providing a heuristic ordering for
candidate generation. Rather than imposing rigid structural constraints
or hard-coding specific modifications, relation scores serve as an
adaptive guidance mechanism. They improve search efficiency by
dynamically deactivating irrelevant repair variables, ensuring that the
solver concentrates on dimensions supported by counterexample evidence.

\subsection{Validation, Admissibility, and Refinement}

Candidate validation serves as the primary correctness guarantee in
MPTA-Repair \cite{clarke1999model}. After a candidate edit set is
generated, the repaired model is verified against the entire property
set $\Phi$. If a candidate resolves the triggering violation but
introduces a failure in another property, it is rejected, indicating
that the current property-trace pool is under-constrained. MPTA-Repair
then extracts the newly observed diagnostic trace and appends it to the
pool, steering the subsequent refinement round.

Candidates that pass full-property regression are further evaluated
against a TARTAR-style admissibility criterion based on functional
equivalence \cite{koelbl2019clock}. This step ensures that clock-bound
modifications do not satisfy safety properties by disabling intended
system behaviors. Rather than enforcing timed-language equivalence,
which would conflict with the goal of eliminating violating timed paths,
admissibility is checked by constructing untimed automata from both the
original and repaired models and comparing their accepted untimed
languages \cite{hopcroft2001automata}. A candidate is rejected if it
adds or removes discrete functional paths, ensuring that only timing
characteristics are changed.

The refinement loop therefore establishes a closed-loop mechanism that
pairs local candidate generation with global validation. Local MaxSMT
formulations propose candidate assignments, full-property regression
ensures complete specification compliance, and functional-equivalence
checking preserves the underlying discrete structure. Blocking
constraints prevent recurrence of failed assignments, ultimately
yielding a repair that satisfies the multi-property specification while
preserving the intended untimed model semantics.
